% Generated draft: requires mathematical and source review.
\documentclass[11pt]{article}
\usepackage[margin=25mm]{geometry}
\usepackage{fontspec}
\setmainfont{Latin Modern Roman}
\usepackage{amsmath,amssymb,mathtools,unicode-math}
\setmathfont{Latin Modern Math}
\usepackage{enumitem,needspace,xurl,hyperref}
\setlength{\parindent}{0pt}
\setlength{\parskip}{5pt}
\setlength{\emergencystretch}{3em}
\newcommand{\N}{\mathbb{N}}
\newcommand{\Z}{\mathbb{Z}}
\newcommand{\Q}{\mathbb{Q}}
\newcommand{\R}{\mathbb{R}}
\newcommand{\C}{\mathbb{C}}
\newcommand{\F}{\mathbb{F}}
\newcommand{\A}{\mathbb{A}}
\newcommand{\PP}{\mathbb{P}}
\newcommand{\E}{\mathbb{E}}
\newcommand{\eps}{\varepsilon}
\newcommand{\dd}{\,\mathrm{d}}
\begin{document}
\section*{UnsolvedMath ID: 700174}
\textbf{Source problem code:} catalogue entry
\subsection{Definitions and mathematical statement}
Let \(\mathcal O\) and \(\mathcal P\) denote, respectively, the classes called \emph{orderable groups} and \emph{pro-orderable groups} in the supplied source. The record does not define these terms and refers to Problem 1.35, whose text is not supplied. In particular, the intended definition of pro-orderability cannot be recovered from this record alone.

Under the interpretation that ``axiomatizable'' means axiomatizable by a set of first-order sentences in the language of groups \(L_{\mathrm{grp}}=\{\cdot,{}^{-1},1\}\), define
\[
\operatorname{Th}(\mathcal P)=\{\varphi:\varphi\text{ is an }L_{\mathrm{grp}}\text{-sentence and }H\models\varphi\text{ for every }H\in\mathcal P\},
\]
and
\[
\mathcal E(\mathcal P)=\{G:G\text{ is a group and }G\models\operatorname{Th}(\mathcal P)\}.
\]
Thus \(\mathcal E(\mathcal P)\) is the smallest first-order axiomatizable class of groups containing \(\mathcal P\), with smallest understood with respect to inclusion.

\textbf{Question (A. I. Kokorin, 1969; Kourovka Notebook Problem 3.20).} With the source's intended definitions, is
\[
\mathcal O=\mathcal E(\mathcal P)?
\]
Equivalently, under the stated first-order interpretation, is it true for every group \(G\) that
\[
G\in\mathcal O\quad\Longleftrightarrow\quad G\models\operatorname{Th}(\mathcal P)?
\]
\textbf{Unresolved qualifications.} This is a conditional formalization: the supplied record does not specify the order-invariance convention for orderable groups, the meaning of pro-orderable groups, or the logical meaning of axiomatizability. The missing reference to Problem 1.35 must be checked before these placeholders can be replaced by complete intrinsic definitions.
\subsection{Short English statement}
Kokorin asks whether the orderable groups are exactly the groups in the smallest axiomatizable class containing all pro-orderable groups. Under a first-order interpretation, this asks whether orderability is equivalent to satisfying every group-language sentence true in every pro-orderable group. The supplied record omits the relevant definitions and the referenced Problem 1.35, so the formulation remains conditional.
\subsection{Sources}
\begin{itemize}
\item Supplied local UnsolvedMath registry record, permanent ID \(700174\), exact source problem code \texttt{catalogue entry}, title \emph{Kourovka Notebook Problem 3.20}; upstream revision \texttt{d371ee917ac9647fbb06cd6fc529deec6f3dda82}. The registry attributes the problem to A. I. Kokorin, proposed in 1969.
\item \emph{The Kourovka Notebook (Unsolved Problems in Group Theory)}, 3rd issue (1969), Problem 3.20; also identified in the supplied record with the 21st issue edition (2026), \texttt{arXiv:1401.0300v48}. The supplied collection description names E. I. Khukhro and V. D. Mazurov as editors. The statement includes the cross-reference ``See 1.35,'' but the text of Problem 1.35 is not supplied.
\item Supplied background and status-review note: an explicitly unverified AI claim gives the negative answer ``No, it does not,'' attributed to Achyuth Jayadevan, \emph{Preprint} (2026), \url{https://zenodo.org/records/22721372}. The status-review note reports that this claim is marked with \(\#\) in the October 2026 version of the 21st issue. This is supplied attribution, not verification of the claim or its appearance in that edition.
\item The supplied registry status and canonical status are both \texttt{open}. Its status review is dated \texttt{2026-10-16} and labels the evidence \texttt{unrefereed\_claim}. This is supplied dated curation, not an independently verified literature assessment; it does not establish the asserted negative answer.
\end{itemize}
\end{document}
